\chapter{Type Theory in Brief}
\label{app:types}

Standard references are~\cite{martinlof1984,hottbook2013,girard1987}.

\section*{Judgments and contexts}

The basic judgments are $\Gamma\vdash A\ \mathsf{type}$, $\Gamma\vdash a:A$, and, for equality, $\Gamma\vdash a\equiv b:A$. A context $\Gamma=(x_1:A_1,x_2:A_2,\dots,x_n:A_n)$ is a list in which each $A_k$ is a type in the preceding context. Contexts support \emph{weakening}, the addition of an unused hypothesis, and \emph{substitution}, the replacement of a variable by a term of its type.

\section*{Dependent sums and products}

The dependent product $\Pi_{x:A}B(x)$ has introduction by $\lambda$-abstraction and elimination by application. The dependent sum $\Sigma_{x:A}B(x)$ has introduction by pairs $(a,b)$ with $b:B(a)$ and elimination by projections. When $B$ does not depend on $x$ these reduce to function and product types.

\section*{Identity types}

For $a,b:A$ the identity type $a=_Ab$ has an introduction rule $\mathsf{refl}_a:a=_Aa$ and an elimination rule, path induction. In the homotopy interpretation, a type is a space, a term is a point, and a term of $a=_Ab$ is a path from $a$ to $b$. A type is a \emph{set} if any two parallel paths are equal.

\section*{Linear logic}

In intuitionistic linear logic a context is split into an unrestricted part $\Gamma$ and a linear part $\Delta$. The linear implication $A\multimap B$ has introduction by a function that consumes exactly one $A$ and elimination by application to a term built from a disjoint part of the linear context. Weakening and contraction are admissible on $\Gamma$ and not on $\Delta$. The multiset semantics used in Chapter~\ref{ch:linear} is the Petri-net reading of the multiplicative fragment with atomic propositions as places.
